\section{Exact local RPPR on supplied supports, paths, and trees}
\label{sec:rppr-terminal}

The block theorem also identifies what is easy and what is hard in the RPPR
problem.  Put
\begin{equation}
 w_D:=D^{1/2}\one,
 \qquad
 c_\rho:=b-\alpha\rho w_D.
 \label{eq:shifted-rppr-demand}
\end{equation}

\begin{theorem}[Shifted Dirichlet KKT certificate]
\label{thm:shifted-dirichlet}
Let $S\subseteq V$, and let
\begin{equation}
 z_S:=Q_{SS}^{-1}c_{\rho,S},
 \qquad z_{S^c}:=0,
 \qquad r(z):=b-Qz.
 \label{eq:shifted-dirichlet-solve}
\end{equation}
If
\begin{equation}
 z_S>0,
 \qquad
 0\leq r_i(z)\leq\alpha\rho\sqrt{d_i}
 \quad\bigl(i\in S^c\cap(\partial S\cup\supp(s))\bigr),
 \label{eq:boundary-kkt-certificate}
\end{equation}
then $z=x^\star(\rho)$.  Moreover,
\begin{equation}
 0\leq D^{-1/2}r(z)\leq\alpha\rho\one,
 \label{eq:rppr-residual-certificate}
\end{equation}
so the same point is an unregularized PageRank approximation with a verified
degree-normalized residual gate at scale $\rho$.
\end{theorem}

\begin{proof}
On $S$, \cref{eq:shifted-dirichlet-solve} gives
$(Qz-b)_i=-\alpha\rho\sqrt{d_i}$.  Since $z_i>0$, this is exactly the active
RPPR KKT equality in \cref{eq:shared-rppr-kkt}.  On $\partial S$,
\cref{eq:boundary-kkt-certificate} is equivalent to
$(Qz-b)_i\in[-\alpha\rho\sqrt{d_i},0]$, the inactive KKT interval.  If
$i\in\operatorname{Ext}(S)\setminus\supp(s)$, then $Q_{iS}=0$ and $b_i=0$,
so $r_i(z)=0$.  The displayed test separately covers seed coordinates outside
$S$.  Thus all KKT conditions hold.  Strong convexity makes the minimizer
unique, proving $z=x^\star(\rho)$.  The same coordinatewise argument gives
\cref{eq:rppr-residual-certificate}.
\end{proof}

\begin{corollary}[Linear terminal rung on a forest support]
\label{cor:forest-rppr-terminal}
If $S=S^\star(\rho)$ is supplied and $G[S]$ is a forest, then the RPPR optimum
and the full boundary KKT certificate can be computed in
\begin{equation}
 O\bigl(\vol(S^\star(\rho))+|\supp(s)|\bigr)
 =O(1/\rho+|\supp(s)|)
 \label{eq:forest-rppr-terminal-work}
\end{equation}
work.  In the single-seed setting this is $O(1/\rho)$.  With a supplied
width-$w$ ordering, the corresponding single-seed bound is
$O((w+1)^2/\rho)$.
\end{corollary}

\begin{proof}
On the exact support, the active KKT equations are precisely the shifted
Dirichlet system.  Apply \cref{thm:forest-resolution} or
\cref{thm:width-resolution}, scan the boundary once, and use the source
support-volume bound $\vol(S^\star(\rho))\leq1/\rho$.  Seed coordinates
outside the support are checked directly from the sparse seed input.
\end{proof}

\begin{lemma}[RPPR support components contain seeds]
\label{lem:support-components}
Every connected component of $G[S^\star(\rho)]$ intersects $\supp(s)$.  In
particular, for a single seed, $S^\star(\rho)$ is either empty or connected
and contains the seed.
\end{lemma}

\begin{proof}
Suppose an active component $C$ contains no seed coordinate.  Because $C$ is
a component of the active induced graph, every active coordinate outside $C$
has zero coupling to $C$, and every inactive coordinate has value zero.  The
active KKT equations on $C$ therefore reduce to
\[
 Q_{CC}x_C^\star=-\alpha\rho D_C^{1/2}\one<0.
\]
The principal matrix $Q_{CC}$ is a positive definite Stieltjes matrix, so
$Q_{CC}^{-1}\geq0$.  Hence $x_C^\star<0$, contradicting positivity on the
support.  For a single seed, two different active components cannot both
intersect its singleton support.
\end{proof}

The lemma rules out disconnected active islands in the single-seed problem.
It justifies boundary expansion as a complete discovery mechanism, although
it does not yet bound how many scans that expansion needs.
